\documentclass[script,pro]{debate}
\motion{我厌蠢，我错了 / 我厌蠢，我没错}
\stance{正方}{我错了}
\meta{赛制}{中科大校赛}
\meta{口径来源}{备赛文档 第十节 正方口径表}
\begin{document}
\debatetitle
\begin{thesisbox}[全场一句话立论]常人的厌蠢，平常凭比较就判人蠢，尺子握在已经会了的人手里，自己看不见；就算看准了，劝了还犯，也会从“说了那么多遍”走到“他就是那么蠢”，回头不再讲理。这份厌平常的样子，该改。\end{thesisbox}

\begin{speech}{1. 正一开篇立论稿}{正文 829 字 / 180 秒}
谢谢主席，各位好。

厌蠢，是觉得别人犯蠢时的嫌恶、反感，常带着想躲开的心；可以只冲着那件事，也可以冲着这个人、看低他；它是态度，不等于骂人或欺负人。蠢，是在一般人看来本该想到、做对的事上明显想错、做错，不论出于不会还是不肯。错没错，看清醒的我问完三句：看准了吗？对他公道吗？划得来吗？这份厌平常的样子该改，就是错了；只在个别场合调调分寸，就是没错。

这里的我，是一个普通的、会说出这句话的我。我方要证明，这份厌平常的样子有该改的毛病；对方要证明，它平常站得住。门槛双方共用：平常就这样、只是改得掉，照样算错了。

我方第一点：拿我当一般人。

常人的厌蠢，平常凭比较下判断，不凭记录。教过、被指出、又做错，这是记录，我们不争。可厌蠢的典型说法是比较。一位博主教人坐高铁，评论区有人留言：坐高铁不是谁都会吗？这个也要教？新京报的文章写同学、同事，写的是“我都可以做到你为什么做不到”的抱怨。

这把尺子，握在已经会了的人手里。1975 年的经典实验发现：知道结果的人，会高估不知道的人本来能知道多少，自己还几乎察觉不到。2023 年的实验里，知道答案的大学生，替不认识的同龄同学估字谜难易，偏差和替自己估基本一样。

你想想，尺子只往“他本该会”那头偏。这不是哪一次眼花，是平常那把尺子。要改的，是凭比较、不问就判。这一点回扣看准了吗，看准是公道的前提。

我方第二点：说了也没用。

就算看准了，熟人劝了还犯，常人的厌平常会从“说了那么多遍”走到“他就是那么蠢”，回头不再讲理。

每日商报 2024 年 1 月写过一位程序员：同事工作两年多，还常填错表格；他多次劝告，同样的错误还是接二连三，关系一度降至冰点，现在很微妙。报道紧接着写了一句：说了那么多遍，怎么还是那么蠢。

恼火我们不剔，要看它往哪走。2007 年的三项研究发现：生气过后会和解，轻蔑才把人推出去，还可能从一次次生气里长出来；觉得管不了他、把错归到他这个人，都预示轻蔑。说白了，“说了那么多遍”是管不了，“还是那么蠢”是归到了人。

越线看回头：这回先不说，不越线；回头他开口先认定不值得听，就越线了。这一点回扣对他公道吗里的尊重。尊重是底线。

我们全场都会请对方回答：\stress{常人说“我厌蠢”的那一刻，说出口的是“你卡在哪”，还是“这都不会”？}

我们不要你第一秒不厌；我们要你问完三句之后，别再认同那份看低人的厌。谢谢。
\end{speech}
\begin{contingency}
\ifthen{对方说“题目写的是厌，不是鄙视”}{两头都写了，都用“可以”，没替任何一方选；平常落在哪一头、落在那里错不错，是今天要辩的。}
\ifthen{对方问“看错了也算厌蠢？”}{算。厌蠢说的是你的厌；看没看准，是第一问要称的，看得准的厌也在里面。}
\ifthen{对方立“对得上对象就没错”的判准}{对得上只回答了一问；它没回答厌了以后还给不给他作为人的尊重，也没回答划得来吗。契合不等于道德上恰当。}
\ifthen{对方说“我有权远离”}{有权这样做，不等于这样做没错；题目问的是这份厌对不对，不是别人能不能管你。}
\ifthen{对方说“情绪不经选择，无所谓对错”}{题目问的是厌，不是骂；“无所谓对错”也不等于“没错”。第一秒的厌我们不评，评的是第二秒以后你认不认同它。}
\end{contingency}

\begin{stage}{2. 正一被质询防守预案}{反二质询，120 秒双边计时}
\textbf{原则}：先答是或否，再用一句话守住前提，每答不超过十秒。反二的取证五问（准、走神、恼火、一时走开、一般人）全队同一答案：不争、不争、不反对、不算、是；答完第五问补一句我方争的两处。被逼二选一时，指出问题的预设，给出我方表述。被问对屡教不改的人怎么办，只答表里那一句：对那件事生气、批评都可以。
\begin{qa}
\qarow{教过、指出过、又做错，这种看准你们争不争？}{不争。凭记录的看准，我们让。}
\qarow{走神、手滑这类犯蠢，你们争不争准？}{不争。}
\qarow{对那件事本身的恼火，你们反不反对？}{不反对。对那件事生气、批评都可以。}
\qarow{这回先不说、一时走开，算不算越线？}{不算。越线按时间划：回头他开口先认定不值得听，合作只下指令不给理由，才越线。}
\qarow{“一般人”是了解他处境的一般人吗？}{是，不是评论区的多数票。我们争的是两处：没有记录、凭比较就判的那一下，和看准之后回头还讲不讲理。}
\qarow{“清醒的我”谁说了算？}{清醒只是说双方摆出来的事实都知道：蠢的成因和代价，厌的后果和功能。你们的事实也在里面。}
\qarow{“划得来吗”，那有用就是对的？}{有用要算，但它是三问之一，不是全部；只赢一问不自动赢下整场，对你们对我们都一样。}
\qarow{他不肯学，厌他就公道了吧？}{公道问两件事：该不该由他负责，厌了以后还给不给他作为人应得的尊重。不肯学，只回答了前一件。}
\qarow{“也可以冲着这个人看低他”，是不是把鄙视塞进了厌？}{两头都写了，都用“可以”，没替任何一方选。平常落在哪一头、落在那里错不错，是今天要辩的。}
\qarow{只在心里厌、从不表现，算不算厌蠢？}{算。心里那份就是厌蠢本身；它错不错，看判准。}
\qarow{调整分寸也是改，岂不是一点不完美就算错？}{不是。只在个别场合冒出来的，调分寸就好；要改的是这份厌平常的样子，才叫错。}
\qarow{你们的实验估的都是陌生人吧？}{是不认识的同龄人。可比较也在同学同事之间：新京报的文章写同学、同事，写的也是“我都可以做到你为什么做不到”。身边不等于有记录。}
\qarow{回头不讲理，报道写了吗？}{没写，这一步是我们的推论。依据：轻蔑不看好对方会变好；认定人改不了，人就少教；康德也问，否认他有理智，又怎能让他明白自己错了。}
\qarow{对屡教不改的人，你们要我们怎么办？}{对那件事生气、批评都可以，一时走开也可以；要改的是看低他，把他划到不用再讲理那边。}
\end{qa}
\end{stage}

\begin{stage}{3. 正二质询反一问题链}{120 秒，双边计时}
\textbf{总目标}：拿到三个承认。一，没有记录时，说他本该会，凭的是比较；常人说出口的是反问。二，看准了不等于可以看低这个人，回头先认定不值得听算越线。三，门槛规则对方也认：平常就这样、只是改得掉，照样算错。对方说超过十五秒还没回到问题，礼貌打断；一条链拿不到就换下一条。
\begin{chain}{链一：没有记录的那一下}{逼出“凭比较”与“反问”}
\q{没有记录时，说他本该会，凭的是不是比较？}{答“是”：下一问；答“不是”：那凭什么？请举一样}
\q{“这都不会？”是在问他，还是已经判完了？}{答“判完了”：下一问；答“在问”：那他答“我不会”，您方厌不厌？}
\q{常人厌蠢时先开口问，您方有材料吗？}{答“没有”：收束；答“有”：请报出处，我方驳论核对}
\cut{好，我理解为您方没有否认。下一个问题。}
\close{谢谢。没有记录的那一下，凭的是比较；常人先问，对方拿不出材料。}
\end{chain}
\begin{chain}{链二：就算看准了}{逼出“看准不等于可以看低人”与越线线}
\q{看准了他这件事蠢，就能看低他这个人吗？}{答“不能”：下一问；答“能”：那尊重小问，您方怎么过？}
\q{回头他开口，先认定不值得听，算不算越线？}{答“算”：下一问；答“不算”：那什么才算越线？}
\q{“怎么还是那么蠢”，说的是事还是人？}{答“人”：收束；答“事”：那“那么蠢”的主语是谁？}
\cut{时间有限，我换个问法：“那么蠢”说的是不是他？}
\close{谢谢。看准不等于可以看低人；回头先不听就越线，而那句话落在人身上。}
\end{chain}
\begin{chain}{链三：平常还是个别}{逼出门槛规则}
\q{平常就这样、只是改得掉，照样算错，您方认吗？}{答“认”：下一问；答“不认”：任何毛病都改得掉，那分界看什么？}
\q{先问、回头讲理，是常人平常就有，还是要学才有？}{答“要学”：收束；答“平常就有”：请举一个常人先问的例子，不是老师，不是圣人}
\cut{好，我理解为您方没有否认这条规则。}
\close{谢谢。要学才有，就是平常没有；平常没有，照样算错。}
\end{chain}
\end{stage}

\begin{kit}{4. 正二驳论}{套件 · 组装后 520–542 字 / 120 秒}
\assembly{开头 → 回应反二（任选 1） → 驳对方立论（任选 2） → 收尾}
\begin{fixed}{开头}
谢谢主席。我先说一句：我方想请评委把对方的立论放到判准下看，问完三句，常人这份厌平常的样子，到底站不站得住。
\end{fixed}
\begin{pick}{回应反二}{1}
\begin{module}{如果反二打的是我方实验估的都是陌生人}
反二说，我方的实验估的都是陌生人。凭记录的那一格，我们本来就不争。可比较不挑场合：新京报那篇文章写同学、同事，写的也是“我都可以做到你为什么做不到”。身边不等于有记录。记录要一次次攒，“这都不会？”那一下，只要比较就够。
\end{module}
\begin{module}{如果反二打的是回头不讲理是推论}
反二说，回头不讲理是推论。这一点我方承认，报道没写。推论的依据有三条：轻蔑不看好对方会变好；认定人改不了，人就少教，经理研究里看法一改，辅导的意愿就上来了；康德也问，在一个判断上否认他有理智，又怎能让他明白自己错了。
\end{module}
\begin{module}{如果反二打的是先问本来就在反方的厌里}
反二说，先问本来就在对方的厌里。那请拿出常人先问的样子。调研到的典型说法，“怎么会有人不会坐飞机啊”，“你能不能有点常识”，全是反问。反问是判完了，询问才是在问。要靠立论承诺才补得上，恰恰说明这份厌自己不往那里走。
\end{module}
\begin{module}{如果反二打的是我方只证明了有，没证明平常}
反二说，我方只证明了有，没证明平常。门槛规则写着：平常就这样、只是改得掉，照样算错了。我们不说每次判错，说的是方向只往一头偏、自己看不见、每回都用。这样的尺子，就是平常。“个别”要证，不能靠喊，对方得拿出常人平常先问的样子。
\end{module}
\begin{fallback}{反二没有打到我方论点}
我方两个论点，分工很清楚。第一点管这份厌凭什么起：没有记录，就凭比较，尺子握在会了的人手里，自己看不见。第二点管它往哪走：就算看准了，劝了还犯，也会从“说了那么多遍”走到“他就是那么蠢”。这两处，都等着对方正面回答。
\end{fallback}
\end{pick}
\begin{pick}{驳对方立论}{2}
\begin{module}{如果对方立论建立在「教过了还犯」上}
对方立论的主干，是教过了还犯。这一格，我方在准上让：教过、被指出、又做错，凭记录的看准我们不争。可它让对方赢到的，只是冲着不肯的恼火公道。对他公道吗还有后一小问：厌了以后，还给不给他作为人应得的尊重。“他不肯学”只回答了前一小问。你想想，劝了还犯，正好把轻蔑的两个预测因素供上：“说了那么多遍”，是觉得管不了他；“还是那么蠢”，是把错归到他这个人。所以教过了还犯，恰恰是我方第二点开始的地方。
\end{module}
\begin{module}{如果对方立论建立在「恼他不弃他」上}
对方立论的主干，是恼他不弃他：恼里带着讲理，所以不是轻蔑。恼火我方本来不剔。可“不纯”不等于“没有”。2007 年那篇研究写着：两种情绪常常一起出现，轻蔑可能从先前的生气里长出来。恼不是挡住轻蔑的盾，是它生长的地方。Strawson 那句也只划了两头：态度完全客观，才可以跟他谈、甚至谈判，却不能跟他讲理。他没说，有恼就一定在讲理。还在恼，只说明还没到完全客观；回头讲不讲理，对方得另证。
\end{module}
\begin{module}{如果对方立论建立在「情绪按能不能控制分流」上}
对方立论的主干，是情绪按能不能控制分流：判断他该负责就生气，控制不了就同情。这条我方全认。可它管的是判断之后生出什么，不管判断本身准不准。毛病出在判断这一步。1988 年的实验发现：被归为心理和行为一类的毛病，人们把起因看作当事人自己能控制的，于是更多愤怒、更少同情；归因信息一换，情绪就跟着变。说白了，没人告诉你他卡在哪，人默认把他读成不肯。厌蠢那一下，恰恰是线索最少的时候。
\end{module}
\begin{module}{如果对方立论建立在「情绪不经选择，无所谓对错」上}
对方立论的主干，是情绪不经选择，无所谓对错。亚里士多德确实说过，我们感到愤怒不经选择；可紧接着他就说，受责备的是以某种方式去感受的人。第一秒的厌，我们不评，双方都认。评的是第二秒以后，常人实际怀着的那份厌：你认同它、放任它，还是修正它。改，是换一种看法，不是压住感受。还有一层，“无所谓对错”不等于“没错”。对方今天要证的是“我没错”，这是一句正面的自我辩护，推定谁也不占，绕不过去。
\end{module}
\begin{module}{如果对方立论建立在「守标准、止损」的好处上}
对方立论的主干，是这份厌守得住标准、止得住损。守标准，对那件事的恼和纠正就守得住，用不着看低人。五个实验、一万多名被试、五十二个议题里，没有一种事实纠正引发逆火：纠正，人是听得进的。可厌一在场，人怕显得无能，就不敢认错、不敢求助。1999 年对一家企业五十一个团队的研究写道：承认错误、求助的人可能显得无能，所以人不敢提。Linux 内核照样拒绝不合格的补丁，守质量，用不着看低人。
\end{module}
\begin{fallback}{对方立论的主干不在以上几条时}
我方想把对方的立论放回判准。对方要证的，不是世上存在一份厌得对的厌，而是常人这份厌平常站得住。只证明这份厌可以停在事上，不够；只证明毛病改得掉，也不够：平常就这样、只是改得掉，照样算错了。存在更好的，不等于这个就错了，这一点我方认，所以我方先证缺陷：平常凭比较，尺子自己看不见；就算看准了，也会从要求走到判词。对方要赢，得正面回答这两处，拿出常人平常的样子。
\end{fallback}
\end{pick}
\begin{fixed}{收尾}
说到底，回到三问。看准了吗：平常凭比较，尺子自己看不见。对他公道吗：就算看准了，劝了还犯，也会从“说了那么多遍”走到“他就是那么蠢”。这份厌平常的样子，该改。
\end{fixed}
\end{kit}

\begin{stage}{5. 正四对辩要点}{90 秒，单边计时}
\textbf{目标}：B 块里把战场钉在两处：没有记录的那一下凭什么，看准之后还讲不讲理。每次发言先一句回应，再抛一个问题；对方连续两次不答同一问，下一次发言点出来：“这个问题对方已经两次没有回答。”
\begin{points}{进攻点（4–5 个，每个 15–20 秒）}
\item 第一次做错的那一下，手里还没有记录。那时说“这都不会？”，凭的是记录，还是比较？
\item 常人说“我厌蠢”的那一刻，说出口的是“你卡在哪”，还是“这都不会”？请拿出常人先问的样子。
\item 看准了他这件事，就能看低他这个人吗？“说了那么多遍，怎么还是那么蠢”，落在事上，还是人上？
\item 知道答案的人自己察觉不到偏差，要求忽略答案也没用。那“我看得很准”，靠什么保证？
\item 平常就这样、只是改得掉，照样算错了。先问、回头讲理，是常人平常就有，还是要学才有？
\end{points}
\begin{points}{备用点（6–8 个）}
\item 对方说熟人之间记录在手：记录要一次次攒；新京报的文章写同学、同事，写的也是“我都可以做到你为什么做不到”。身边不等于有记录。
\item 对方说情绪按能不能控制分流：那条只管判断之后，不管判断准不准；没人告诉你他卡在哪时，人默认把他读成不肯。
\item 对方说恼还在就不是轻蔑：两种情绪常常一起出现，轻蔑可能从一次次生气里长出来；恼不是盾，是它生长的地方。
\item 对方说一时走开不越线：我们同意。越线看回头：他开口先认定不值得听，合作只下指令不给理由。
\item 对方说存在更好的不等于这个就错了：我们认，所以先证缺陷，再说换成什么。
\item 对方举孔子骂了子路照样讲完：那是圣人、是老师，正是我们要常人学的改法，不是常人平常的样子。
\item 对方说后见之明平均才偏一截：我们不说偏得大；方向只往一头偏，自己看不见；同一篇还写，设法减少偏差的研究，效应量没有更低。
\item 对方说研究要被试设想同伴经历和自己相同：这条局限我们认；另一项没有这条指示，判断自己的知识在同龄人里有多普遍，熟悉感就足以把判断带偏。
\end{points}
\begin{points}{对方最可能问的三个问题 → 回应}
\item 教过了还犯，还不能怪？→ 能怪，怪那件事；要改的是“怪”走到看低他、回头不再跟他讲理那一步。
\item 回头不讲理，报道写了吗？→ 没写，是我们的推论。依据是轻蔑不看好对方会变好；认定人改不了，人就少教；康德也问，否认他有理智，又怎能让他明白自己错了。
\item 你们的实验估的都是陌生人？→ 凭记录的我们本来不争；比较也在同学同事之间，身边不等于有记录。
\end{points}
\end{stage}

\begin{stage}{6. 正三盘问问题链}{90 秒，单边计时}
\textbf{原则}：对方开始展开就打断：“谢谢，我听到了答案。”链一同一问先问反二、再问反四，两人不一致就是小结的第一句。
\begin{chain}{链一：先问还是先判}{问二辩 → 问四辩同一问题，检验口径}
\q{（问反二）没有记录时，常人是先问，还是先判？}{答“先问”：转问反四；答“先判”：转问反四，看是否一致}
\q{（问反四）同一个问题：没有记录时，常人是先问，还是先判？}{与反二一致：下一问；不一致：记下，小结重锤}
\q{（问反四）先问，是常人平常就有，还是要学才有？}{答“要学”：记下；答“平常就有”：有没有一个常人先问的例子？}
\end{chain}
\begin{chain}{链二：看准之后}{钉越线线与那句判词}
\q{（问反二）回头他开口先认定不值得听，算不算越线？}{答“算”：下一问；答“不算”：那什么才算？}
\q{（问反四）“怎么还是那么蠢”，说的是事还是人？}{答“人”：下一问；答“事”：那“那么蠢”的主语是谁？}
\q{（问反四）下回还讲理，是常人自然会，还是要提醒自己？}{答“要提醒”：收束；答“自然会”：记下，自由辩追问材料}
\close{谢谢两位。线是对方划的，那句话落在人身上。}
\end{chain}
\textbf{盘问目标}：小结里证明两件事。一，没有记录时，对方说不出常人先问的样子，或者两位辩手答得不一样；二，回头先认定不值得听是越线，而“那么蠢”落在人身上。两件事分别接回我方第一点“看准了吗”和第二点“对他公道吗”里的尊重。
\end{stage}

\begin{stage}{7. 正二、正四被盘问防守预案}{两人口径必须一致}
\textbf{原则}：回答短；不确定时说“我方立论的立场是……”；两人的回答要能互相引用。被问是不是推论，直接认，再给依据。
\begin{qaa}
\qaarow{你们让了准，论点一是不是只剩网上陌生人？}{不是。比较也在同学同事之间：新京报的文章写同学、同事，写的是“我都可以做到你为什么做不到”。身边不等于有记录。}{不是。身边不等于有记录；新京报的文章写同学、同事，写的也是“我都可以做到你为什么做不到”。}
\qaarow{你们说的平常，除了评论区和一则报道，还有什么？}{平常靠机制：比较那一下随手可用，记录要攒；知道答案的人自己察觉不到偏差。方向单一、看不见、每回都用的尺子，就是平常。}{平常靠机制：记录要攒，比较随手可用；知道答案的人自己察觉不到。方向单一、看不见、每回都用，就是平常。}
\qaarow{回头不讲理，是不是推论？}{是推论，报道没写。依据：轻蔑不看好对方会变好；认定人改不了，人就少教；康德问，否认他有理智，又怎能让他明白自己错了。}{是推论。依据是轻蔑不看好对方会变好；看法一改，经理辅导的意愿就上来；康德那句附释。}
\qaarow{恼火你们不反对，去掉看低后剩下的那份，错没错？}{剩下的对那件事的恼，我们不反对。可要改的那一截看低，是这份厌平常的样子，所以这份厌本身该改，哪怕改完还剩一点反感。}{对事的恼我们不反对；看低是平常的样子，所以这份厌本身该改，哪怕改完还剩一点反感。}
\qaarow{对一个共事两年、劝过多次的同事，“先开口问”问什么？}{问他卡在哪。问了、他说不会又不学，归责就到你们那边了；我们要改的，只是不问就判，和回头不再讲理。}{问他卡在哪。问完他不学，归责在他；我们要改的是不问就判，和回头不再讲理。}
\qaarow{存在更好的，就说这个错了？}{不是。存在更好的，不等于这个就错了，这一点我们认。所以我们先证缺陷，再说换成什么。}{不是，这一点我们认；先证缺陷，再说换成什么。}
\qaarow{你们是不是要我们别生气、做圣人？}{不是。恼照样恼，对那件事照样批评，一时走开也行；只是把“这都不会？”换成“你卡在哪？”，回头照样跟他讲理。}{不是。改是换一种看法，不是压住感受；恼照样恼，回头照样讲理。}
\qaarow{研究让被试设想同伴经历和自己相同，是不是推着人投射？}{这条局限我们认。另一项没有这条指示：判断自己的知识在同龄人里有多普遍，熟悉感就足以把判断带偏。}{局限我们认；另一项研究没有这条指示，熟悉感照样把判断往“大家都知道”那头带。}
\qaarow{只在心里厌，不表现，也错？}{心里那份就是厌蠢本身，它错不错看判准；心里看低他，问完三句同样站不住。}{心里那份就是厌蠢本身，错不错看判准；看低他，问完三句同样站不住。}
\qaarow{你们说的“我”，是不是网上的喷子？}{不是。一个普通的、会说出这句话的我；不是网暴者，也不是圣人。}{不是。一个普通的、会说出这句话的我。}
\end{qaa}
\end{stage}

\begin{kit}{8. 正三盘问小结}{套件 · 组装后 519–528 字 / 120 秒}
\assembly{开头 → 链一的结果（任选 1） → 链二的结果（任选 1） → 接回我方论点}
\begin{fixed}{开头}
谢谢主席。我先说清楚：我方的盘问围着两件事，一件是这份厌凭什么起，一件是它往哪走。我一件一件说给评委听。
\end{fixed}
\begin{pick}{链一的结果}{1}
\begin{module}{如果对方说没有记录时常人也是先问}
对方说，没有记录的时候，常人也是先问。好，请评委记住这句。再看调研到的典型说法：“怎么会有人不会坐飞机啊？我厌蠢症犯了。”新京报写到的“你能不能有点常识”。“怎么还有人到医院不会挂号呀？太蠢了吧。”全是反问。反问是判完了，询问才是在问。如果先问真是常人平常的样子，说出口的不会全是反问。对方要拿出常人先问的样子，不是老师，也不是圣人。
\end{module}
\begin{module}{如果对方承认没有记录时常人是先判}
对方承认了：没有记录的时候，常人是先判。那判凭的是什么？没有记录，只剩比较。而比较这把尺子，握在已经会了的人手里。知道结果的人，会高估不知道的人本来能知道多少，自己还几乎察觉不到；替一个不认识的同龄人估，也一样偏。所以我方第一点那一格，对方自己替我们证完了：没有记录的那一下，凭的就是比较，尺子只往“他本该会”那头偏。
\end{module}
\begin{module}{如果反二、反四答得不一样}
同一个问题，问了对方两位辩手，给出的答案不一样。这件事值得评委记下来。没有记录时，常人是先问还是先判，正是对方要守住“平常”的那一格；守这一格的两个人，自己没有对齐。对方要证这份厌平常站得住，至少得先说清，它平常是怎么下判断的。我方的答案一直是一句：没有记录，就凭比较，尺子握在会了的人手里，只往“他本该会”那头偏，自己看不见。
\end{module}
\begin{module}{如果对方回避了链一的问题}
没有记录时，常人是先问还是先判，这个问题，对方没有正面回答。请评委记住。为什么难答？因为调研到的典型说法，全是反问：“怎么会有人不会坐飞机啊？”“你能不能有点常识？”反问是判完了，询问才是在问。说先问，拿不出常人的样子；说先判，就认了凭比较。我方的答案很简单：没有记录，就凭比较。对方要证这份厌平常站得住，这一格绕不过去，迟早要答。
\end{module}
\begin{fallback}{链一没有问完}
第一条链要问的，我替评委说完：没有记录的时候，说他本该会，依据只剩比较。比较这把尺子，握在已经会了的人手里，只往“他本该会”那头偏；知道结果的人会高估不知道的人本来能知道多少，自己还几乎察觉不到。所以要改的，是凭比较、不问就判：没有记录，先开口问他的处境。对方要证这份厌平常站得住，就得拿出常人先问的样子来。
\end{fallback}
\end{pick}
\begin{pick}{链二的结果}{1}
\begin{module}{如果对方承认回头先认定不值得听算越线}
对方承认了：回头他开口先认定不值得听，就越线。那请回到每日商报那一案：劝了多次，关系一度降至冰点，现在很微妙，报道紧接着写下：说了那么多遍，怎么还是那么蠢。“说了那么多遍”，是觉得管不了他；“还是那么蠢”，是把错归到他这个人。这两样都预示轻蔑，而轻蔑不看好对方会变好。不看好他会变好，回头还听他的理由吗？线是对方自己划的。
\end{module}
\begin{module}{如果对方说「那么蠢」说的是事}
对方说，“怎么还是那么蠢”说的是事。那请问，“那么蠢”的主语是谁？是他。心理学说，轻蔑把对方贬为愚蠢、无能，贬的是品性多于行为。“蠢”字一落下来，就落在人身上。劝的时候，说的是一件事：工作上要多上心。到了这一句，说的是这个人：那么蠢。说白了，从要求走到判词，就是从对事走到对人。这正是我方第二点说的那一步，对方的回答替我们走完了。
\end{module}
\begin{module}{如果对方回避了链二的问题}
回头他开口先认定不值得听，算不算越线；“那么蠢”说的是事还是人。这两个问题，对方没有正面回答，请评委记住。这条线本来是双方共用的：这回先不说、一时走开，不越线；回头先认定不值得听，越线。对方不在这里表态，就说不清这份厌平常停在线的哪一边。何况劝了还犯之后，说出口的那句“还是那么蠢”，已经把错归到了他这个人，落在人身上了。
\end{module}
\begin{fallback}{链二没有问完}
第二条链要问的，我替评委说完：就算看准了，劝了还犯的时候，“说了那么多遍”是觉得管不了他，“还是那么蠢”是把错归到他这个人，这两样都预示轻蔑；轻蔑可能就从一次次的生气里长出来。线双方共用：这回先不说，不越线；回头他开口先认定不值得听，就越线。从要求走到判词，就是从对事走到对人。常人的厌平常停在哪一边，对方要拿出样子。
\end{fallback}
\end{pick}
\begin{fixed}{接回我方论点}
所以盘问指向的，是我方两个论点，也是判准的两问。

第一点管这份厌凭什么起，回扣看准了吗。凭记录的我们让；没有记录的那一下，凭的是比较，尺子握在会了的人手里，只往一头偏，自己看不见。第二点管它往哪走，回扣对他公道吗里的尊重。就算看准了，劝了还犯，也会从要求走到判词。恼火我们不剔，要改的是看低他、回头不再讲理那一截。

你想想，门槛规则写着：平常就这样、只是改得掉，照样算错了。对方要赢，就得证常人平常有记录、平常先问、平常回头讲理。这三样，请评委在自由辩里替我们盯着。
\end{fixed}
\end{kit}

\begin{stage}{9. 自由辩战场设计}{240 秒}
\begin{battleground}{战场一：“这都不会？”凭的是什么（前 90 秒，主战场，回扣看准了吗）}
\begin{fields}
\field{目标}{让评委看到：没有记录的那一下，常人凭比较下判断，说出口的是反问；凭比较的尺子握在会了的人手里，自己看不见。}
\field{开场问题}{没有记录时，说他本该会，凭的是什么？}
\field{对方答“熟人有记录” → 追问}{第一次做错那一下，记录在哪？记录要一次次攒，“这都不会？”只要比较就够。新京报的文章写同学、同事，写的也是“我都可以做到你为什么做不到”。}
\field{对方答“凭一般人的标准” → 追问}{定义里的一般人了解他的处境；“我都可以做到”量的是我，不是一般人。知道答案的人自己察觉不到偏差，要求忽略答案也没用。}
\field{对方可能反攻 → 我方防守}{“你们的实验估的都是陌生人。”→ 是不认识的同龄人；凭记录的我们本来不争，身边不等于有记录。“平均才偏一截。”→ 我们不说偏得大，方向只往一头偏、自己看不见；设法减少偏差的研究，效应量没有更低。}
\field{收束语}{没有记录的那一下，凭的就是比较；尺子在会了的人手里，自己看不见。}
\end{fields}
\end{battleground}
\begin{battleground}{战场二：“说了那么多遍”之后（中间 90 秒，回扣对他公道吗）}
\begin{fields}
\field{目标}{让评委看到：就算看准了，劝了还犯，常人的厌会从要求走到判词，回头不再讲理；越线线双方共用。}
\field{开场问题}{看准了他这件事蠢，就能看低他这个人吗？}
\field{对方答“不能” → 追问}{那“说了那么多遍，怎么还是那么蠢”，落在事上还是人上？“说了那么多遍”是管不了，“还是那么蠢”是归到人，两样都预示轻蔑。}
\field{对方答“看低能力可以” → 追问}{区别评价可以，但必须受承认尊重约束；回头他开口先认定不值得听，算不算越线？}
\field{对方可能反攻 → 我方防守}{“恼还在，就不是轻蔑。”→ 两种情绪常常一起出现，轻蔑可能从一次次生气里长出来；还在恼只说明没到完全客观。“回头不讲理是推论。”→ 认；依据是轻蔑不看好对方会变好，认定人改不了人就少教。}
\field{收束语}{怪那件事，我们不拦；要改的是从要求走到判词、回头不再讲理那一截。}
\end{fields}
\end{battleground}
\begin{battleground}{战场三：平常还是个别（最后 60 秒前插入，回扣门槛）}
\begin{fields}
\field{目标}{把对方的“个别场合”逼成要证的事：常人平常有记录、平常先问、平常回头讲理。}
\field{开场问题}{平常就这样、只是改得掉，算不算错？}
\field{对方答“算” → 追问}{那先问、回头讲理，是常人平常就有，还是要学才有？}
\field{对方答“不算” → 追问}{任何毛病都改得掉，那分界看什么？双方立论共用的规则，您方现在改口吗？}
\field{对方可能反攻 → 我方防守}{“存在更好的，不等于这个就错了。”→ 我们认，所以先证缺陷再说替代。“孔子骂了子路照样讲完。”→ 圣人、老师，正是我们要常人学的改法。}
\field{收束语}{要学才有，就是平常没有；平常就这样、只是改得掉，照样算错了。}
\end{fields}
\end{battleground}
\begin{points}{必问（前两分钟内问完）}
\item 常人说“我厌蠢”的那一刻，说出口的是“你卡在哪”，还是“这都不会”？（全场问题）
\item 没有记录时，说他本该会，凭的是不是比较？
\item 看准了他这件事，就能看低他这个人吗？
\item 平常就这样、只是改得掉，算不算错？
\end{points}
\begin{points}{必答（15 秒以内正面回答）}
\item 对屡教不改的怎么办？→ 对那件事生气、批评都可以，一时走开也可以；要改的是看低他，把他划到不用再讲理那边。
\item 只在心里厌，算不算？→ 算。心里那份就是厌蠢本身；它错不错，看判准。
\item 回头不讲理是不是推论？→ 是，报道没写；依据是轻蔑不看好对方会变好，认定人改不了，人就少教。
\item 你们是不是要我们别生气？→ 不是。改是换一种看法，不是压住感受；恼照样恼，回头照样讲理。
\end{points}
\textbf{时间分配与站位}：前 90 秒打战场一，一辩开场问、二辩接追问（二辩驳论打过的点由二辩接）；中间 90 秒打战场二，三辩带盘问成果上、四辩补越线线；剩约 60 秒时插一轮战场三，由四辩收束，重复“平常就这样、只是改得掉，照样算错了”，不开新战场；留 20 到 30 秒备用。对方开新战场，与胜负无关的不接，一句话带回战场一。一辩守定义与判准，二辩守论点一，三辩守论点二，四辩负责收束与价值。
\end{stage}

\begin{kit}{10. 正四结辩}{套件 · 组装后 892–920 字 / 210 秒}
\assembly{归纳分歧（任选 1） → 处理最强点（任选 1） → 回应反四（任选 1） → 临场 → 论点收束 → 价值升华 → 结尾}
\begin{pick}{归纳分歧}{1}
\begin{module}{如果全场主战场落在「平常凭比较还是凭记录」}
谢谢主席。今天全场打得最多的，是平常凭比较还是凭记录。双方其实有共识：教过、被指出、又做错，凭记录的看准，我们不争。分歧只在一格。没有记录的那一下，常人凭的是什么？我方说，记录要一次次攒，“这都不会？”那一下，只要比较就够。凭记录的那一格，从来不是分歧。评委要比的，是常人平常落在哪一格。
\end{module}
\begin{module}{如果全场主战场落在「看准之后还讲不讲理」}
谢谢主席。今天全场打得最多的，是看准之后还讲不讲理。线按时间划：这回先不说、一时走开，不越线；回头他开口先认定不值得听，越线。分歧只在常人平常落在线的哪一边。我方拿出的，是劝了还犯之后说出口的那一句：怎么还是那么蠢。那一句里，要求已经变成了判词。就算看准了，也会走到这一步。
\end{module}
\begin{module}{如果全场主战场落在「先问是不是常人平常就有」}
谢谢主席。今天全场打得最多的，是先问是不是常人平常就有。没有记录就该先问，这是我方要的改法。分歧在于，这是常人平常的样子，还是要改出来的样子。我方说，调研到的典型说法全是反问，“怎么会有人不会坐飞机啊”。反问是判完了，询问才是在问。说本来就有，就得拿出常人先问的样子。
\end{module}
\begin{fallback}{主战场不在以上几条时}
谢谢主席。全场打到最后，其实落在一处：常人说“我厌蠢”时，那份厌平常停在事上，还是会从“这都不会？”走到“他就是那么蠢”，回头不再跟他讲理。判准里有一条规则：平常就这样、只是改得掉，照样算错了。要比的，是谁描述的平常，更像一个普通的、会说出这句话的人。看准了吗、对他公道吗，都落在这里。
\end{fallback}
\end{pick}
\begin{pick}{处理对方最强点}{1}
\begin{module}{如果对方全场最有力的是「熟人之间记录在手」}
对方全场最有力的一点，是熟人之间记录在手。我方承认，有记录的那一格，看得准，我们让。可身边不等于有记录。新京报那篇文章写同学、同事，写的是“你能不能有点常识”的烦躁，是“我都可以做到你为什么做不到”的抱怨。记录要一次次攒；第一次做错的那一下，手里什么也没有，只有比较。说白了，凭比较的那一下不挑场合。而且就算记录在手，也只过了看准这一关；看准之后往哪走，正是我方第二点。
\end{module}
\begin{module}{如果对方全场最有力的是「回头不讲理是推论」}
对方全场最有力的一点，是回头不讲理是推论。这一点我方承认：报道没写，这一步双方都只能推。可推论有依据。轻蔑不看好对方会变好；2006 年的三项研究发现，经理怎么看人能不能改，预示他后来辅导下属的多少，看法一改，辅导的意愿和建议的质量都上来了。康德也问：在一个判断上否认他有理智，又怎能让他明白自己错了？我方不说从此不理，报道写的现状是很微妙，不是照常教、照常听。
\end{module}
\begin{module}{如果对方全场最有力的是「恼里带着讲理，不是轻蔑」}
对方全场最有力的一点，是恼里带着讲理，所以不是轻蔑。恼火我方不剔，这个立场没变过。可“不纯”不等于“没有”。2007 年的研究写得很清楚：两种情绪常常一起出现，轻蔑可能就从对同一个人的一次次生气里长出来。恼不是挡住轻蔑的盾，是它生长的地方。还在恼，只说明还没到完全客观；回头讲不讲理，要另证。“怎么还是那么蠢”这一句里，要求已经变成了判词。要求说的是事，判词说的是人。
\end{module}
\begin{module}{如果对方全场最有力的是「存在更好的，不等于这个就错了」}
对方全场最有力的一点，是存在更好的，不等于这个就错了。这一点我方认。所以我方没有拿“有更好的态度”去证错，而是先证缺陷：平常凭比较，尺子只往一头偏，自己看不见；就算看准了，也会从要求走到判词。缺陷证在前，替代才登场。何况平常就这样、只是改得掉，照样算错了。替代是换一种看法，不是压住感受：恼照样恼，对那件事照样批评，只是先开口问一句你卡在哪，回头照样跟他讲理。
\end{module}
\begin{fallback}{对方最有力的一点不在以上几条时}
我方想请评委再听一遍我方最硬的一条。1975 年的经典实验发现：知道结果的人，会高估不知道的人本来能知道多少，自己还几乎察觉不到。2023 年的实验里，知道字谜答案的大学生，替一个不认识的同龄同学估难易，偏差和替自己估基本一样；要求他们忽略答案，也没用。看不见，也压不下。这样一把尺子，常人说“这都不会？”的时候，每回都在用。说白了，这不是哪一次眼花，是平常那把尺子；看准，又是公道的前提。
\end{fallback}
\end{pick}
\begin{pick}{回应反四}{1}
\begin{module}{如果反四收在「你本可以」的恼是看得起他}
反四最后说，恼里那句“你本可以”，是看得起他。我方同意，要求是一种尊重，所以对事批评，我们一直支持。问题出在要求落空之后：说了那么多遍，要求就变成了判词，怎么还是那么蠢。报道里劝了多次之后，紧接着写下的就是这一句。从要求走到判词，就从交涉走到了退出。看得起他的，是前半句。我们要改的，是后半句。
\end{module}
\begin{module}{如果反四收在「只在个别场合冒出来，那叫调分寸」}
反四最后说，误判和过头只在个别场合，调调分寸就好。门槛规则写着：平常就这样、只是改得掉，照样算错了。个别要证，不能靠喊。我方证在平常的有两处：一把方向固定、自己看不见的尺子；一条劝了还犯就把轻蔑的预测因素供上的路。对方要说个别，得拿出常人平常先问、平常回头讲理的样子。
\end{module}
\begin{module}{如果反四收在「存在更好的，不等于这个就错了」}
反四最后说，存在更好的，不等于这个就错了。这句话我方认。可我方从来没有拿“有更好的”去证错。我方证的是缺陷：平常凭比较，尺子看不见；就算看准了，也会从要求走到判词。缺陷在前，替代在后。替代登场时，它也不难，是换一种看法，不是压住感受：恼照样恼，把“这都不会？”换成“你卡在哪？”，回头照样讲理。
\end{module}
\begin{module}{如果反四收在「先问、讲理本来就在我们的厌里」}
反四最后说，先问、讲理，本来就在他们的厌里。那这正是我方要的改法。可判准评的是常人那份厌平常的样子，不是反思修好的版本。调研到的典型说法全是反问。说本来就有，就得拿出常人先问的样子。如果先问要靠立论承诺才补得上，恰恰说明这份厌自己不往那里走。平常就这样、只是改得掉，照样算错了。
\end{module}
\begin{fallback}{反四的收束不在以上几条时}
反四的结辩，请评委用判准去量：问完三句，常人这份厌平常的样子，站不站得住。看准了吗：平常凭比较，尺子握在会了的人手里，自己看不见。对他公道吗：就算看准了，劝了还犯，也会从要求走到判词，回头不再讲理。按门槛规则，对方要说个别，得拿出常人平常的样子。还有一句，我直接回应。
\end{fallback}
\end{pick}
\live{60}
\begin{fixed}{论点收束}
我方的两个论点，打到现在，还站着。

第一点，拿我当一般人，回扣看准了吗。凭记录的看准，我们让；没有记录的那一下，常人凭的是比较，同学同事之间也一样。尺子握在已经会了的人手里，只往“他本该会”那头偏，自己还看不见。这不是哪一次眼花，是平常那把尺子。

第二点，说了也没用，回扣对他公道吗里的尊重。就算看准了，劝了还犯，“说了那么多遍”是管不了，“还是那么蠢”是归到人。说白了，恼火我们不剔；要改的，是从要求走到判词、回头不再讲理的那一截。你想想，这两处，哪一处是个别场合？
\end{fixed}
\begin{fixed}{价值升华}
我想请评委想一个很小的时刻。

周五下午五点，办公室里，同事第三次把表格递过来，又填错了。他站在你桌边，手里还捏着那支笔。你接过来，叹了口气。那一秒的烦，我们不评，谁都会有，我也有。

接下来是第二秒。你可以说：“这都不会？”也可以说：“你卡在哪？”

两句话，只差几个字。说出口，不到两秒。前一句说完，他就成了一个不用再解释的人。后一句说完，他还是一个可以讲理的人。

我们今天争的，不是你该不该烦，是烦完之后，\stress{你还愿不愿意听他把话说完}。

你我也都会有递过表格、等着对面叹气的那一天。那一天，你希望对面那个人，说的是哪一句？
\end{fixed}
\begin{fixed}{结尾}
恼，照样可以恼；只是别把一个人，判到不用再讲理的那一边。所以，我厌蠢，我错了。谢谢。
\end{fixed}
\end{kit}
\begin{contingency}
\ifthen{时间不够}{先删处理最强点模块里的研究年份和细节，再把回应反四压成第一句；论点收束、价值升华与结尾不删。}
\ifthen{反四在结辩里提出全场没出现过的新论据}{临场位只用一句点出“这是结辩里的新论据，全场没有经过检验”，不展开，直接进论点收束。}
\ifthen{反方在自由辩之后发动奇袭质询，追问“平常还有什么材料”}{正四在结辩临场位里正面答：平常靠机制，比较随手可用、记录要攒，尺子看不见；回头不讲理是推论，依据是轻蔑不看好对方会变好。}
\end{contingency}

\begin{stage}{11. 正方奇袭预案}{质询 150 秒双边 / 申论 120 秒，二选一}
\textbf{执行人}：奇袭质询由正三执行，奇袭申论由正一执行；由执行人在环节结束时举手向主席示意发动。不在开场就用掉，越晚越清楚对方哪里最疼；进入反四结辩就没有机会了。
\begin{fields}
\field{窗口一：反一立论之后}{发动条件：反一把蠢收窄成“明知故犯”“又菜又犟”，或把厌写成“只讨厌蠢事”。发动奇袭质询，用链一钉定义。执行人：正三。反一用的是中立定义，不发动。}
\field{窗口二：对辩之后}{发动条件：反四在对辩里说出“看准了就可以看低他”这类与尊重小问相抵的话。发动奇袭质询，用链三取证。执行人：正三。}
\field{窗口三：盘问小结之后、自由辩之前}{发动条件：我方论点一在盘问小结里被打成“只剩陌生人”，或“先问是不是常人平常就有”还没有人正面打。发动奇袭申论，立“反问与询问”战场，给自由辩定议题。执行人：正一。}
\field{窗口四：自由辩之后、结辩之前}{发动条件：对方全场把“平常”压在熟人有记录上，或自由辩里说“先问、回头讲理是常人平常的样子”却没有给出常人的材料。发动奇袭质询，用链二或链四，逼反四在结辩里回答。执行人：正三。首选窗口。对方没有可以钉死的表态时，改用奇袭申论，执行人正一。}
\end{fields}
\begin{chain}{奇袭质询链一：如果对方把蠢收窄成明知故犯、又菜又犟}{逼出“辩护的是厌坏”}
\q{只厌又菜又犟的人，那是厌蠢，还是厌坏？}{答“厌蠢”：下一问；答“厌坏”：收束}
\q{对一个真不会、在学的人，厌他算不算错？}{答“错”：下一问；答“没错”：那您方定义就没剔掉不会}
\q{当场分得清他是不想知道，还是不能知道吗？}{答“分不清”：收束；答“分得清”：靠记录还是靠比较？}
\cut{好，我理解为您方没有否认。}
\close{谢谢。对方辩护的，一半是厌坏；剩下那一半，当场分不清。}
\end{chain}
\begin{chain}{奇袭质询链二：如果对方说先问、回头讲理是常人平常的样子}{逼出“要学才有”}
\q{常人厌蠢时先开口问，您方有一条材料吗？}{答“没有”：下一问；答“有”：请报出处}
\q{“这都不会？”是在问他，还是已经判完了？}{答“判完了”：下一问；答“在问”：那他答“我不会”，您方厌不厌？}
\q{要提醒自己才会先问，算不算平常就有？}{答“不算”：收束；答“算”：那提醒从哪来？}
\close{谢谢。先问要靠提醒，就是平常没有。}
\end{chain}
\begin{chain}{奇袭质询链三：如果反四在对辩里说看准了就可以看低他}{钉尊重小问}
\q{看准了他这件事，就可以看低他这个人吗？}{答“可以”：下一问；答“不可以”：那您方四辩刚才那句收回吗？}
\q{看低他这个人，还守得住作为人应得的尊重吗？}{答“守得住”：下一问；答“守不住”：收束}
\q{回头他开口先认定不值得听，算不算越线？}{答“算”：收束；答“不算”：那什么才算？}
\close{谢谢。看低到人，尊重就守不住；线是对方自己认的。}
\end{chain}
\begin{chain}{奇袭质询链四：如果对方全场把平常压在熟人有记录上}{逼出“记录要攒，比较随手”}
\q{熟人第一次做错，您手里有记录吗？}{答“没有”：下一问；答“有”：第一次的记录是什么？}
\q{没记录那一下，“这都不会？”凭的是不是比较？}{答“是”：下一问；答“不是”：那凭什么？}
\q{记录要攒，比较随手就有，平常先用的是记录，还是比较？}{两种答法都收束到“第一次那一下凭的是比较”}
\close{谢谢。记录要一次次攒，第一次那一下，凭的就是比较。}
\end{chain}

\begin{kit}{奇袭申论}{套件 · 组装后 520–523 字 / 120 秒}
\assembly{开头 → 为什么现在发动（任选 1） → 立一个战场 → 收尾}
\begin{fixed}{开头}
谢谢主席。我方用掉这一次奇袭，只做一件事：请评委看一个我方准备好、立论里还没展开的战场：反问和询问。
\end{fixed}
\begin{pick}{为什么现在发动}{1}
\begin{module}{如果在盘问小结之后发动}
盘问小结刚刚结束，自由辩还没开始。我方想先替自由辩定一个议题，请双方在接下来的自由辩里正面打：常人说“我厌蠢”的那一刻，说出口的是询问，还是反问。我方先把答案摆出来。这个问题，决定双方说的“平常”，哪一边更像常人。
\end{module}
\begin{module}{如果在自由辩之后发动}
自由辩刚刚结束，结辩还没开始。我方想在最后两篇发言之前，把一个问题留给对方四辩，请他在结辩里正面回答：常人说“我厌蠢”的那一刻，说出口的是询问，还是反问。我方先把答案摆出来。这个问题，决定双方说的“平常”，哪一边更像常人。
\end{module}
\begin{fallback}{发动时机不在以上几个窗口时}
我方选在这里发动，是因为有一个问题，双方都要在剩下的环节里回答：常人说“我厌蠢”的那一刻，说出口的是询问，还是反问。双方都绕不开它。我先把我方的回答摆出来，请对方接着答。这个问题，决定双方说的“平常”，哪一边更像常人。
\end{fallback}
\end{pick}
\begin{fixed}{立一个战场}
先看说出口的话。调研到的厌蠢典型说法：“怎么会有人不会坐飞机啊？我厌蠢症犯了。”“怎么还有人到医院不会挂号呀？太蠢了吧。”新京报那篇写同学、同事，是“你能不能有点常识”。你听，全是反问。

反问和询问，差的不是语气。询问是还没判，在等他的答案；反问是判完了，答案已经在问的人心里。

那判凭什么？1988 年的实验发现：被归为心理和行为一类的毛病，人们把起因看作当事人自己能控制的，于是更多愤怒、更少同情；归因信息一换，情绪就跟着变。说白了，没人告诉你他卡在哪的时候，人默认把他读成不肯。你想想，厌蠢那一下，正是线索最少、判得最快的时候。

所以我方要的改法，就落在一个问号上。后见之明的综述说，鼓励人去考虑别的原因，能减少这种偏差。把别的原因问出来，最直接的办法，就是开口问一句：你卡在哪？要提醒自己才会先问，就说明平常不先问。

恼照样恼。只是先问，再判。
\end{fixed}
\begin{fixed}{收尾}
所以回到判准的第一问：看准了吗。没有记录，先问，才看得准；平常就这样、只是改得掉，照样算错了。请对方回答：常人说“我厌蠢”的那一刻，说出口的，是询问，还是反问？
\end{fixed}
\end{kit}
\textbf{不发动的情况}：场上没有出现可以钉死的表态，只是我方预设的矛盾；我方两个论点都没有被打塌，而且自由辩的战场已经在盘问小结里立起来；对方已经在自由辩之后发动奇袭质询，正四要把时间留给结辩里正面回答；反四结辩已经开始。宁可不用，也不要把奇袭申论念成第二篇立论。
\end{stage}
\end{document}
